\textbf{Solução:}

\begin{itemize}
    \item \textbf{(a) 5,2} (O 3º algarismo é 3, então mantém-se o 2).
    \item \textbf{(b) 0,83} (Os primeiros algarismos significativos são 8 e 2. O 3º é 5, logo arredonda-se para cima).
    \item \textbf{(c) 81} (O 3º algarismo é 3, então mantém-se o 1).
    \item \textbf{(d) 3,1} ($\pi \approx 3{,}14159\ldots$, o 3º algarismo é 4, então mantém-se o 1).
    \item \textbf{(e) 6,2} (O 3º algarismo é 3, então mantém-se o 2).
    \item \textbf{(f) 0,043} (Zeros à esquerda não contam. O 3º algarismo é 8, logo arredonda-se o 2 para 3).
    \item \textbf{(g) 10} (O 3º algarismo é 4, então mantém-se o zero).
    \item \textbf{(h) 4,0} (O 3º algarismo é 3, então mantém-se o zero, que aqui é significativo).
\end{itemize}